\section*{序論}
10月31日朝までの生成AI関連Xでは、モデルの大規模化と公開形態、サイバー防御、推論インフラが並行して話題になった。Mistral Large 4のオープンウェイト公開とGemini 4 Carbonの限定提供が比較される一方、GoogleのProject Suncatcherは宇宙のTPUを用いた応答生成を紹介した。新しいモデルの能力だけでなく、実際にどこで、どの条件で動かすかが焦点になっている。

運用と社会実装では、OpenAIの監視体制・Model Spec、カナダ政府によるCohere支援、Anthropicの科学研究支援が議論された。さらに音声エージェント企業の買収やデザイン向けモデル、検索経由のトラフィック変化も取り上げられ、生成AIの影響が開発環境から産業政策、オンラインサービスの事業構造まで広がっている。

本号はGrokが収集した12件を掲載順に記事化した。主要なモデル発表の一部は観測窓より前で、窓内の反応と区別する。数値・買収・利用条件など公式一次資料を直接確認できていない項目は、投稿上の主張として記述する。独立した追加調査は行っていない。
